% !TeX root = compte-rendu.tex
% !TeX spellcheck = fr


\section{Introduction}

Le parlement européen est un organisme dont toutes les décisions ont énormément d’influence. Il est donc nécessaire que ces décisions soient prises de manière équitable. Or, ce n’est aujourd’hui pas le cas. Notre objectif va être de définir et modéliser ce que nous entendons par « équitable » et de proposer le système de vote qui répond mathématiquement à ces modèles.

Tout d’abord, nous poserons les fondements de notre travail en nous appuyant sur la théorie des jeux et la celle des indices de pouvoir. Ensuite, nous appliquerons cette théorie au système de vote du parlement européen pour en démontrer les inégalités. Enfin, nous vous présenterons ce qui serait selon nous le meilleur système de vote que le conseil de l’UE pourrait adopter.